\section{Components, forms, and rational maps along a path}
\label{lift:section}

We now work in the remaining situation of Proposition~\ref{grw:remaining-paths}.
In particular the characteristic is fixed, the classical ranks tend to infinity,
and the geometric steps with a nontrivial unipotent radical have been excluded.
All statements about classical groups below are made beyond the finitely many
small parameters already discarded in that proposition.  We write
$P^{(\infty)}$ for the terminal member of the derived series of a finite group
$P$.  A group is \emph{quasisimple} if it is perfect and its quotient by its
centre is nonabelian simple.  Soluble groups can occur beside an active
component; their presence must be distinguished from a noncentral soluble
extension underneath it.

We need two kinds of structure from the remaining finite comparisons.
The normal-subgroup tests will be applied to actual quasisimple components
and their centralizers.  Later, the irreducible-triple classification will
be applied to connected algebraic groups containing the finite images.
Proposition~\ref{lift:edge-lifts} supplies both the components and rational
lifts of their finite inclusions.  Lemma~\ref{lift:path-lifts} composes
these lifts along one chosen path: irreducibility on its finite image
then implies irreducibility on the containing algebraic image.
Independently chosen rational paths need not agree, even when their
finite inclusions do.

The preparatory lemmas extract components from the retained subgroup
structures, choose standard finite covers uniformly over the field,
and lift the prescribed quadratic form in characteristic two.

\subsection{Extracting actual finite components}

\begin{lemma}[A two-step component lemma]\label{lift:component-preimage}
Suppose $P_0\leq P_1$, $E_1\lhd P_1$ is quasisimple, and $P_1/E_1$ is
soluble.  Let $\pi:P_1\longrightarrow\operatorname{Aut}(E_1)$ be conjugation.
Suppose that $\overline P_0=\pi(P_0)$ has a normal quasisimple subgroup
$\overline E_0$ with soluble quotient.  Then
\[
 E_0=P_0^{(\infty)}\leq E_1
\]
is quasisimple and $P_0/E_0$ is soluble.  In particular $E_0$ is normalized
by every group normalizing $P_0$.
\end{lemma}

\begin{proof}
Set $Q=P_0\cap E_1$ and $A=\pi(Q)$.  The group $\overline P_0/A$ is a
quotient of $P_0/Q$, which embeds in $P_1/E_1$, so it is soluble.  A perfect
group has trivial image in a soluble group.  Consequently
$\overline E_0\leq A$.  Let $B$ be its full inverse image in $Q$.  Then
$B\lhd P_0$ and
\[
 B/Z_0\cong\overline E_0,
 \qquad Z_0=B\cap\ker\pi=B\cap Z(E_1)\leq Z(B).
\]
Because $B/Z_0$ is perfect, $B=B'Z_0$, and therefore
$B'=[B'Z_0,B'Z_0]=B''$.  Thus $B'$ is perfect.  The inverse image in $B$
of $Z(\overline E_0)$ has central commutators, and its intersection with
$B'$ is central in $B'$: for an element of this inverse image, commutation
with $B'$ is a homomorphism from the perfect group $B'$ to the abelian
group $Z_0$.  It is trivial.  It follows that $B'$ is a perfect central
extension of the nonabelian simple group
$\overline E_0/Z(\overline E_0)$, hence is quasisimple.

The quotients $P_0/Q$, $Q/B\cong A/\overline E_0$, and $B/B'$ are
soluble.  Hence $P_0/B'$ is soluble.  Since $B'$ is perfect, it equals
$P_0^{(\infty)}$.  Finally $B'\leq Q\leq E_1$, and the terminal derived
subgroup is characteristic.
\end{proof}

\begin{remark}\label{lift:not-composition-count}
The hypothesis about $\overline E_0$ concerns its structure as a normal
subgroup.  The assertion that $\overline P_0$ has just one nonabelian
composition factor would not suffice: a semidirect product of a nontrivial
module by a simple group can have that property without having a component.
The explicit structures in Proposition~\ref{cls:catalogue} supply the
stronger hypothesis used here.
\end{remark}

\subsection{Standard finite covers at every field size}

\begin{lemma}[Standard finite covers]\label{lift:finite-covers}
There is an absolute rank bound, independent of the finite field,
with the following property.  Let $\mathbf G$ be a simply connected
classical algebraic group above this bound, let $\sigma$ be a standard
split or graph-twisted Steinberg map, and put $U=\mathbf G^\sigma$.
Then $U$ is perfect, $H_2(U,\mathbb Z)=0$, and its standard simple
quotient map $u:U\twoheadrightarrow S$ is central.  For any quasisimple
central extension $a:L\twoheadrightarrow S$, there is a unique
homomorphism $\pi:U\to L$ satisfying $a\pi=u$.  It is surjective and
has central kernel.  A specified projective representation of $L$
also lifts uniquely to a linear representation of $U$ over an
algebraically closed field.
\end{lemma}

\begin{proof}
Steinberg's multiplier Theorem has only a finite list of exceptions
for the split groups, and the accompanying twisted statement has
another finite list
\cite[Theorem~1.1 and the opening paragraph of p.~528]{Steinberg1981}.
Their ranks are bounded, independently of the field size.  We supply
perfectness separately: the root-presentation statement
\cite[Lemma~$32'$, original pp.~51--53]{SteinbergNotes} applies to the
universal split groups, and
\cite[Theorem~34 and its proof, original pp.~187--188]{SteinbergNotes}
gives the twisted case outside its finite list of exceptions.
The root-generated groups used there are the full fixed-point groups
for simply connected groups over finite fields: see
\cite[Corollary~3 to Theorem~7, original p.~65, and the Corollary
after Lemma~64, original p.~186]{SteinbergNotes}.
Increase the discarded rank bound beyond these finite exception lists.
Thus $H_1(U,\mathbb Z)=H_2(U,\mathbb Z)=0$ and $u$ is the usual
central simple quotient.  This argument retains every field size,
including $q=2,3,4$.

Let $A=\ker a\leq Z(L)$.  Pull back $a$ along $u$ to obtain the
central extension
\[
 1\longrightarrow A\longrightarrow U\times_S L
       \longrightarrow U\longrightarrow1.
\]
The action on $A$ is trivial.  The universal coefficient sequence
\[
 0\longrightarrow
 \operatorname{Ext}^1_{\mathbb Z}(H_1(U,\mathbb Z),A)
 \longrightarrow H^2(U,A)
 \longrightarrow\operatorname{Hom}(H_2(U,\mathbb Z),A)
 \longrightarrow0
\]
gives $H^2(U,A)=0$.  The extension splits, and projection of a
splitting to $L$ supplies $\pi$.  Two such maps differ by a
homomorphism $U\to A$, which is zero because $U$ is perfect.
If $H=\pi(U)$, then $L=HA$ and $A$ is central.  Therefore
\[
 L=[L,L]=[HA,HA]=[H,H]\leq H,
\]
so $\pi$ is onto.  Its kernel is contained in $\ker u\leq Z(U)$.

For a projective representation, compose with $\pi$ and pull back
$\operatorname{GL}(V)\to\operatorname{PGL}(V)$.  This is a central
extension of $U$ by $k^\times$.  The same cohomology argument splits
it; perfectness gives uniqueness of the resulting linear lift.
\end{proof}

\subsection{Quadratic forms in characteristic two}

Here and below a Frobenius twist $M^{[a]}$ means pullback by the $p^a$-power
Frobenius morphism.  In particular $[1]$ denotes the first Frobenius twist;
we write $[f]$ when $q=p^f$.  For a dominant weight $\lambda$, the symbols $L(\lambda)$,
$\nabla(\lambda)$, and $\Delta(\lambda)$ denote the simple,
costandard, and Weyl rational modules of highest weight $\lambda$,
respectively; we use the highest-weight, simple-head and simple-socle
properties in \cite[Part~II, Chapter~2]{Jantzen2003}.  Write $X^+$ for the dominant weights and $w_0$ for
the longest Weyl-group element, and put $\lambda^*=-w_0\lambda$.

\begin{lemma}[Restriction of first cohomology]\label{lift:h1-restriction}
Let $\mathbf G$ be a simply connected simple algebraic group over an
algebraically closed field $k$ of characteristic $p$.  Let
$\sigma=F_q\gamma$, where $q=p^f\geq2$ and $\gamma$ is a standard diagram
automorphism commuting with $F_q$.  For every nonzero $q$-restricted
dominant weight $\nu$, restriction is injective:
\[
 H^1(\mathbf G,L(\nu))\longrightarrow
 H^1(\mathbf G^\sigma,L(\nu)).
\]
\end{lemma}

\begin{proof}
Put $I=\operatorname{ind}_{\mathbf G^\sigma}^{\mathbf G}k$.
The induction filtration of
\cite[arXiv version, Proposition~3.1.2]{BNPPSS2015}, obtained from the
coordinate-ring filtration and the Lang map, gives $I$ a filtration with
sections
\[
 \nabla(\lambda)\otimes\nabla(\lambda^*)^\sigma,
 \qquad \lambda\in X^+,
\]
each once.  Write such a section as
$\nabla(\lambda)\otimes\nabla(\mu)^{[f]}$, where $\mu$ is obtained
from $\lambda^*$ by the diagram automorphism.  Choose the regular height
\[
 h(\xi)=\left\langle\xi,\sum_{\alpha>0}\alpha^\vee\right\rangle.
\]
It is invariant under diagram automorphisms and duality, is positive on
every nonzero dominant weight and every positive root, and satisfies
$h(\mu)=h(\lambda)$.

We first justify removing the actual constants from this filtration.
A nonconstant section cannot contain a trivial submodule.  Indeed a
map from $k$ into it would transpose to a nonzero map
$\Delta(\lambda^*)\to\nabla(\mu)^{[f]}$.  The image contains the
simple socle $L(q\mu)$, so its highest weight must occur in the
source.  This gives $qh(\lambda)\leq h(\lambda)$, impossible when
$\lambda\neq0$.  Let $I_j$ be the first filtration stage containing
the actual one-dimensional space of constants.  That space injects
into $I_j/I_{j-1}$, so this section is the unique section for
$\lambda=0$.  The constants split it and
$I_j=I_{j-1}\oplus k$.  The images of the earlier stages in $I/k$,
followed by $I_i/k$ for $i\geq j$, give a filtration of $I/k$;
omitting the repeated stage removes exactly the constant section.
Thus its sections are precisely the nonconstant ones.

We prove that a nonzero $q$-restricted simple $L(\nu)$ cannot map to a
nonconstant section.  Suppose otherwise.  Transposing the second tensor
factor gives a nonzero homomorphism
\[
 L(\nu)\otimes\Delta(\mu^*)^{[f]}
       \longrightarrow\nabla(\lambda).
\]
Its image contains the simple socle $L(\lambda)$.  Filter the Weyl module
by simple composition factors and apply Steinberg's tensor product
Theorem~\cite[Theorems~5.1 and~6.1]{Steinberg1963}.  Since $\nu$ is $q$-restricted, each simple
composition factor of the source has highest weight $\nu+q\eta$ for
some dominant $\eta$.  Consequently
\begin{equation}\label{lift:height-lower}
 \lambda=\nu+q\eta,\qquad h(\lambda)\geq h(\nu).
\end{equation}
Transposing the first tensor factor instead gives a nonzero homomorphism
\begin{equation}\label{lift:transpose}
 \Delta(\lambda^*)\otimes L(\nu)
       \longrightarrow\nabla(\mu)^{[f]}.
\end{equation}
Its image contains $L(q\mu)$, so the weight $q\mu$ occurs in its source.
The maximum height of a weight of that source is
$h(\lambda^*)+h(\nu)=h(\lambda)+h(\nu)$.  Thus
\begin{equation}\label{lift:height-upper}
 qh(\lambda)\leq h(\lambda)+h(\nu).
\end{equation}
Equations~\eqref{lift:height-lower} and~\eqref{lift:height-upper} are
inconsistent when $q>2$.

When $q=2$, both inequalities are equalities.  Hence $\eta=0$ and
$\lambda=\nu$.  The only weight of maximum height in the source of
\eqref{lift:transpose} is $\lambda^*+\nu$, and its weight space is the
one-dimensional span of the tensor $a\otimes b$ of highest-weight
vectors.  If $2\mu\neq\lambda^*+\nu$, the required weight does not
occur.  Otherwise the homomorphism must be nonzero on $a\otimes b$.
Choose a simple root $\alpha$ with
$\langle\nu,\alpha^\vee\rangle=1$, which is possible because $\nu$ is
nonzero and $2$-restricted.  For the Chevalley raising and lowering
operators,
\[
 e_\alpha(a\otimes f_\alpha b)=a\otimes b.
\]
The target of~\eqref{lift:transpose} is Frobenius twisted, so its Lie
algebra action is zero.  Every module homomorphism to it kills the
displayed vector, a contradiction.

A homomorphism from a finite-dimensional simple module to $I/k$ has
image in a finite stage of the filtration.  Taking the first stage with
nonzero intersection with that simple image gives a nonzero map into
one of its sections.  We have proved
\[
 \operatorname{Hom}_{\mathbf G}(L(\nu),I/k)=0.
\]
Apply this also to the dual of the coefficient module.  Tensoring
$0\to k\to I\to I/k\to0$ with $L(\nu)$ shows that the map from
$H^1(\mathbf G,L(\nu))$ to $H^1(\mathbf G,L(\nu)\otimes I)$ is
injective.  The tensor identity for induction and derived Frobenius
reciprocity identify this with the asserted restriction map.  For split
groups this is also the special case $U=k$ of
\cite[Theorem~7.4]{CPSvdK1977}; the argument above includes the diagram
twist and the equality case $q=2$.
\end{proof}

\begin{proposition}[Lifting a prescribed quadratic form]
\label{lift:quadratic}
Let $k$ have characteristic two and let $\mathbf G$, $q$, and $\sigma$
be as in Lemma~\ref{lift:h1-restriction}.  Suppose that $V=L(\nu)$ is
nontrivial and $q$-restricted, and that its restriction to
$\mathbf G^\sigma$ is irreducible and preserves a nonzero quadratic
form $Q$.  If $V$ is not the natural module for a group isomorphic to
$\operatorname{Sp}_{2n}$, $n\geq1$, then
the rational action of $\mathbf G$ preserves the same quadratic form
$Q$.
\end{proposition}

\begin{proof}
The polar form $B$ of $Q$ is nonzero: otherwise $Q$ is the square of a
linear form, and that linear form would be invariant under the
nontrivial irreducible finite action.  Its radical is invariant, so $B$
is nondegenerate.  Thus the finite representation is self-dual.  The
injective parametrization by $q$-restricted highest weights in
Steinberg's finite-group restriction Theorem~\cite[Theorem~43]{SteinbergNotes}
implies that $V\cong V^*$ rationally.  Schur's Lemma, first for
$\mathbf G$ and then for $\mathbf G^\sigma$, shows that their spaces of
invariant bilinear forms are both one-dimensional and agree.
In particular $B$ itself is $\mathbf G$-invariant.

Polarization gives an exact sequence of rational modules
\begin{equation}\label{lift:polarization}
 0\longrightarrow(V^*)^{[1]}
 \longrightarrow\operatorname{Sym}^2(V^*)
 \longrightarrow\bigwedge^2(V^*)\longrightarrow0.
\end{equation}
The first map sends a twisted linear form to its square.  The
obstruction to refining $B$ by a rationally invariant quadratic form
is its connecting class
\[
 c_B\in H^1(\mathbf G,(V^*)^{[1]}).
\]
This class restricts to zero on $\mathbf G^\sigma$, because $Q$ is a
refinement there.  This description is also
\cite[Proposition~3.6 and Corollary~3.9]{GaribaldiNakano2016}.

Let $\mathbf G_1$ be the first Frobenius kernel.  Its action on
$(V^*)^{[1]}$ is trivial.  The inflation--restriction sequence therefore
begins
\[
 0\longrightarrow H^1(\mathbf G/\mathbf G_1,(V^*)^{[1]})
 \longrightarrow H^1(\mathbf G,(V^*)^{[1]})
 \longrightarrow
 \bigl(H^1(\mathbf G_1,k)\otimes(V^*)^{[1]}\bigr)
       ^{\mathbf G/\mathbf G_1}.
\]
The required first-kernel calculation is
\[
 H^1(\mathbf G_1,k)^{[-1]}\cong
 \begin{cases}
 N,&\mathbf G\cong\operatorname{Sp}_{2n},\ n\geq1,
       \quad N\text{ its natural module},\\
 0,&\text{otherwise}.
 \end{cases}
\]
It is reproduced in the proof of
\cite[Proposition~7.1, p.~411]{GaribaldiNakano2016}, citing
\cite[Theorem~3.1(C)(a),(f)]{BNP2004}.  Here the symplectic case includes
$A_1=C_1$ and $B_2=C_2$.  The exceptional module is $L(\omega_1)$ in
$C_n$ numbering, including $n=1$, and $L(\omega_2)$ in $B_2$
numbering.  It is simple and self-dual.  After untwisting, the last invariant space
is therefore zero under our hypothesis.  Inflation, followed by
untwisting the quotient by $\mathbf G_1$, gives an isomorphism
\[
 H^1(\mathbf G,V^*)\xrightarrow{\ \sim\ }
 H^1(\mathbf G,(V^*)^{[1]}).
\]

Write $c_B=F^*c$, where $F$ is ordinary Frobenius.  Put
$\Gamma=\mathbf G^\sigma$ and $M=V^*$.  Since $F$ commutes with
$\sigma$, its restriction $F_\Gamma$ is an automorphism of $\Gamma$.
Naturality gives the commutative square
\[
\begin{array}{ccc}
 H^1(\mathbf G,M)&\xrightarrow{\ F^*\ }&H^1(\mathbf G,M^{[1]})\\
 \operatorname{res}_M\downarrow&&\downarrow\operatorname{res}_{M^{[1]}}\\
 H^1(\Gamma,M)&\xrightarrow{\ F_\Gamma^*\ }&
 H^1(\Gamma,F_\Gamma^*M).
\end{array}
\]
The lower arrow is an isomorphism, with inverse obtained by pulling
back along $F_\Gamma^{-1}$ with the transported coefficients.
This does not identify $M$ with its Frobenius twist.  Hence the
vanishing of $\operatorname{res}(c_B)$ implies that of
$\operatorname{res}(c)$.  This and
Lemma~\ref{lift:h1-restriction}, applied to $V^*$, imply $c=0$ and
$c_B=0$.  Sequence~\eqref{lift:polarization} supplies a
$\mathbf G$-invariant quadratic refinement $Q_0$ of $B$.

Finally $Q-Q_0$ has zero polarization, hence is the square of a linear
form $\ell$.  Its finite invariance implies that $\ell$ is invariant
under $\mathbf G^\sigma$.  Such a linear form is zero on this
nontrivial irreducible module.  Therefore $Q_0=Q$.
\end{proof}

\begin{remark}\label{lift:natural-exception}
The exclusion of the natural symplectic module, including the
$A_1=C_1$ and $B_2=C_2$ conventions just stated, is required in the
inflation argument.  The geometric natural modules and the special
isogenies in characteristic two are constructed as classical algebraic
maps below; Proposition~\ref{lift:quadratic} is used for the remaining
defining-characteristic nonlinear steps.  Suzuki and Ree endomorphisms
have bounded rank and do not occur here.
\end{remark}

\subsection{The retained edge maps}

\begin{lemma}[A nonsingular-line stabilizer]\label{lift:nonsingular-line}
Let $(W,Q)$ be a nonsingular quadratic space of even dimension $2r+2$
over $\mathbb F_q$, where $r\geq3$ and $q$ is even, and let $v\in W$ satisfy
$Q(v)\neq0$.  Write $B$ for the polar form and $\overline W=v^\perp/\langle
v\rangle$.  The kernel of
\[
 O(W,Q)_{\langle v\rangle}\longrightarrow
 \operatorname{Sp}(\overline W,B)
\]
is the central group generated by the symmetry
\[
 r_v(x)=x+\frac{B(x,v)}{Q(v)}v.
\]
The image is the full symplectic group.  The symmetry has nonzero
Dickson invariant, and restriction gives
\[
 \Omega(W,Q)_{\langle v\rangle}
       \cong\operatorname{Sp}_{2r}(q).
\]
In particular this stabilizer has no nontrivial finite unipotent
kernel on $\overline W$.  Over an algebraic closure, the vector
stabilizer $\operatorname{SO}(W,Q)_v$ is a smooth connected adjoint
group of type $B_r$.  The ambient module $W$ restricted to this group has composition factors
$k,N_{B_r},k$ along the flag
$\langle v\rangle<v^\perp<W$.
\end{lemma}

\begin{proof}
A line stabilizer fixes $v$, because a scalar preserving its nonzero
quadratic value has square one, and the field has characteristic two.
Choose a complement $U$ to $\langle v\rangle$ in $v^\perp$ and choose
$t$ with $B(v,t)=1$ and $B(t,U)=0$.  The space $U$ is symplectic.  An
isometry inducing the identity on $\overline W$ sends
$u\in U$ to $u+\ell(u)v$.  Preservation of $Q$ forces $\ell(u)^2Q(v)=0$,
so it fixes $U$ pointwise.  Preservation of $B$ then forces its image
of $t$ to be $t+av$.  Preservation of $Q$ says
$a^2Q(v)+a=0$, hence $a=0$ or $a=Q(v)^{-1}$.  These are exactly the
identity and $r_v$.

For completeness, let $s\in\operatorname{Sp}(U)$.  The difference
$Q(u)+Q(su)$ has zero polarization, so, because the field is perfect,
there is a unique linear form $\ell$ satisfying
\[
 \ell(u)^2Q(v)=Q(u)+Q(su).
\]
The map $v\mapsto v$, $u\mapsto su+\ell(u)v$ is an isometry of
$v^\perp$.  Witt's extension theorem \cite[Theorem~7.4, pp.~57--58]{Taylor1992}
extends this isometry to $W$: the ambient polar radical is zero, so its
radical-intersection criterion is automatic even though $v^\perp$ has
degenerate polar form.
This proves surjectivity.  The formula for $r_v$ also shows that it
commutes with every isometry fixing $v$.

The Dickson invariant of an isometry in even characteristic is
$\operatorname{rank}(g-1)$ modulo two; in particular it is one on
$r_v$.  Its zero kernel is $\Omega(W,Q)$ over a finite field of even
order \cite[Theorems~11.43 and~11.51, pp.~160, 164--165]{Taylor1992};
the exceptional group $\Omega^+(4,2)$ is outside the dimension range here.  Each fibre of the displayed surjection has two elements
differing by $r_v$, so exactly one lies in this kernel.  The last
finite assertion follows.

We also justify the algebraic assertion, including its scheme
convention.  Put $W_0=v^\perp$.  Restriction defines
\[
 H=\operatorname{SO}(W,Q)_v\longrightarrow
       \operatorname{SO}(W_0,Q|_{W_0}).
\]
The kernel of restriction from the full orthogonal vector stabilizer
is the constant group of order two computed above: over an arbitrary
$k$-algebra its remaining parameter satisfies $a+a^2Q(v)=0$, an
\'{e}tale equation.  Its nonidentity element has Dickson invariant one,
so the displayed map has trivial scheme-theoretic kernel.
We give its inverse over an arbitrary $k$-algebra $A$.
Put $c=Q(v)$ and retain the decomposition
$W=kv\oplus U\oplus kt$ used above.  An element
$h\in\operatorname{SO}(W_0)(A)$ has symplectic quotient action, of
determinant one; the odd special orthogonal determinant condition
therefore makes it fix $v$.  Write
$h(u)=su+\ell(u)v$ for $u\in U_A$.  There is a unique $w\in U_A$
such that $B(su,w)=\ell(u)$ for all $u$.  The extension fixing $v$
and restricting to $h$ is obtained by setting
\[
 g(t)=t+w+av,\qquad ca^2+a+Q(w)=0.
\]
The bilinear condition is $B(su,w)+\ell(u)=0$, and the quadratic
condition is exactly the displayed equation.  Its derivative is one
and its leading coefficient is a unit, so adjoining a solution gives
an \'{e}tale double cover.  The two extensions differ by $r_v$ and
have opposite Dickson invariants.  The unique Dickson-zero extension
on this cover agrees on overlaps and descends to $A$.  This is a
functorial inverse, proving the asserted scheme isomorphism.
The odd special orthogonal group is smooth, connected,
and adjoint of type $B_r$ by
\cite[author version, Theorem~C.2.11 and Proposition~C.3.10]{Conrad2014}.
This restriction calculation is also the one in the proof of
\cite[author version, Proposition~C.3.1]{Conrad2014}.
Its action on $W_0/\langle v\rangle$ is the purely inseparable
isogeny to the symplectic group described in
\cite[author version, Remark~C.3.6]{Conrad2014}; consequently this
quotient is the simple module $N_{B_r}$.  Both outer factors of the
flag are trivial.  We used the vector stabilizer: the line-stabilizer
scheme $\operatorname{SO}(W,Q)_{\langle v\rangle}$ is
$H\times\mu_2$.  Indeed its scalar on $v$ has square one,
and multiplying by the inverse central scalar makes it fix $v$;
the two factors have trivial intersection.  Its reduced subgroup
and its finite-field points are the ones just described.
\end{proof}

\begin{proposition}[Components and rational edge maps]
\label{lift:edge-lifts}
In the setting of Proposition~\ref{grw:remaining-paths}, each retained
projected subgroup $P$ has a characteristic normal quasisimple subgroup
$L=P^{(\infty)}$ and $P/L$ is soluble.  In a fixed higher coordinate,
these components nest with the projected subgroups and are normalized
by their common bottom action.  All intermediate components of a
shortcut are normalized by its lower coordinate stabilizer.  The
transported product of an active component maps onto the lower
nonabelian socle and is not contained in the common bottom.

For every retained geometric or defining-characteristic almost-simple
edge, its actual inclusion on these components has a rational lift
between their simple classical algebraic groups, after passage to
simply connected central covers.  The lift agrees with the finite
inclusion after the specified central quotients.

On a geometric edge, every nontrivial absolute composition factor of
the target natural module is a tensor product of Frobenius twists of
the source natural simple module or its dual.  The number of tensor
factors is uniformly bounded by the constants in
Proposition~\ref{grw:remaining-paths}.  Trivial factors and repeated
copies are allowed.  For type $B$ in characteristic two, ``natural
simple module'' means the $2r$-dimensional quotient of the usual
$(2r+1)$-dimensional module.
\end{proposition}

\begin{proof}
\emph{Finite components and their normalization.}
We first specify the components in the structures of
Proposition~\ref{cls:catalogue}.  A nondegenerate block or form subgroup
contains the derived classical group on each active block.  A tensor
subgroup contains the central product of the derived groups on its
tensor factors.  An extension-field or subfield subgroup contains the
derived classical group over its indicated field.  In the retained
twisted tensor construction, the single Frobenius orbit of the tensor
factors has fixed-point group equal to the classical group over the
corresponding extension field; its derived group is the active
quasisimple group.  The almost-simple case contains its quasisimple
socle preimage.

The reductions leave one nonabelian active factor.  Hence every
permitted normalizer preserves it.  All other base factors commute
with it modulo scalars, and their surviving images are soluble.
The retained diagonal, field, graph, and factor-permutation quotients
are soluble as well; the exclusion of other nonabelian factors here
includes factors in the permutation quotient, not just base factors.
Thus the explicitly displayed active derived group is normal and the
quotient is soluble.  Scalar preimages introduce only central kernels:
if $C/Z$ is the displayed perfect active factor, with $Z$ central,
then $C=C'Z$ and $C'=C''$, giving its quasisimple derived preimage.
For the characteristic-two nonsingular-line case,
Lemma~\ref{lift:nonsingular-line} gives the asserted component directly.
For a two-edge projection, we verify the precise hypothesis of
Lemma~\ref{lift:component-preimage}.  Use the common upper-coordinate
notation $P_0\leq P_1$ and $Q_1\leq R_1\leq P_1$ from
Lemma~\ref{cls:core-gaps}.  Here $Q_1$ is soluble, $R_1/Q_1$ is the
middle simple active factor, and $P_1/R_1$ is soluble.  The upper
retained edge supplies the quasisimple component
$E_1=P_1^{(\infty)}$.  Perfectness gives $E_1\leq R_1$.
Its image in $R_1/Q_1$ is nontrivial, since $Q_1$ is soluble, and is
normal, hence is the full simple group.  Therefore
\[
 E_1Q_1=R_1,\qquad E_1\cap Q_1=Z(E_1).
\]
Both $E_1$ and $Q_1$ are normal in $P_1$.  Thus, for $x\in Q_1$,
commutation with $x$ maps $E_1$ into $E_1\cap Q_1=Z(E_1)$.
It is a homomorphism to an abelian group, so perfectness makes it
trivial.  Consequently $Q_1$ centralizes $E_1$.  Conjugation on
$E_1$ agrees with conjugation on
$E_1/Z(E_1)\cong R_1/Q_1$: an automorphism of the perfect group
$E_1$ acting identically modulo its centre differs from the identity
by a homomorphism $E_1\to Z(E_1)$ and is therefore the identity.

The original-core product identity in
Proposition~\ref{frm:active-transport} gives $Q_1\leq P_0$.
Indeed its transported product comes from the middle original core,
which is contained in $D$ and hence in the lower labelled group.
It follows that
\[
 J=(P_0\cap R_1)/Q_1
\]
is the actual first-edge inner projection of
Lemma~\ref{cls:core-gaps}.  Its retained structure supplies the
normal quasisimple component $J^{(\infty)}$ with soluble quotient.
This is a structural assertion about this particular $J$; the
excluded parabolic or module semidirect extensions have not been
inferred to split.

Let $\pi:P_1\to\operatorname{Aut}(E_1)$ be conjugation.  By the
centralization and simple-quotient identification just proved,
$\pi(P_0\cap R_1)$ is the inner image of $J$, isomorphic to $J$.
It is normal in $\pi(P_0)$, and
$\pi(P_0)/\pi(P_0\cap R_1)$ is soluble because it is a quotient of
$P_0/(P_0\cap R_1)\leq P_1/R_1$.
Hence the image of $J^{(\infty)}$ is normal quasisimple in
$\pi(P_0)$, with soluble quotient.  These are exactly the hypotheses
of Lemma~\ref{lift:component-preimage}, which produces the actual
lower component inside $E_1$.

Finally, if $P\leq P'$ and $L'$ is the
component of $P'$, the image of the perfect group $L$ in the soluble
group $P'/L'$ is trivial.  Hence $L\leq L'$.  Characteristicity gives
normalization by the common bottom.  Each intermediate inner
projection is normal in its own coordinate stabilizer, which contains
the lower coordinate stabilizer; its characteristic component is
therefore normalized by that lower stabilizer too.

To check that the active component is not killed, use the notation
$Q\leq R\leq P$ of Proposition~\ref{frm:active-transport}.  In our
multiplicity-one coordinate, $R/Q$ is the nonabelian simple active
factor and $P/R$ is soluble.  Hence $L\leq R$.  Since $P/L$ is
soluble, the simple group $R/Q$ has trivial image modulo $LQ$, and
$LQ=R$.  Thus $L$ maps onto $R/Q$.  Transporting $L$ over all upper
coordinates gives a normal subgroup of the lower labelled group,
mapping onto its socle modulo its original bottom core.  If this
product were contained in the common bottom, normality would put it
inside that core, a contradiction.

For each actual finite component $L_i$, fix its standard cover
$U_i=\mathbf G_i^{\sigma_i}$ and the surjection
$\pi_i:U_i\twoheadrightarrow L_i$ of
Lemma~\ref{lift:finite-covers}, over the specified simple quotient.

\smallskip\noindent
\emph{Rational maps for the geometric edges.}
We have obtained the finite components, their nesting, and their survival
under the original bottom-core quotient.  We now lift their inclusions.
For the geometric cases we describe the maps on natural modules over $k$,
which also proves the assertion about their composition factors.
Write $N$ for the source natural simple module.
\begin{enumerate}[label=\textup{(\roman*)}]
\item A block embedding acts on its active block by $g$, on a dual
block, when present, by $g^{-\mathsf T}$, and trivially on omitted
spaces.  Its constituents are $N$, $N^*$, and trivial modules.
A form embedding is the natural action with its prescribed invariant
bilinear or quadratic form.  These are rational matrix maps.

\item For a tensor factor with an omitted space $U$, the map is
$g\mapsto g\otimes1_U$.  On the active factor its natural module is
the direct sum of $\dim U$ copies of $N$, or of a prescribed dual or
Frobenius twist of $N$.  The commuting omitted factor changes the
multiplicity space, not this formula.

\item Restriction of scalars from $\mathbb F_{p^{ab}}$ to
$\mathbb F_{p^a}$, followed by extension of scalars to $k$, has the
explicit decomposition
\[
 k\otimes_{\mathbb F_{p^a}}\mathbb F_{p^{ab}}
       \cong\bigoplus_{j=0}^{b-1}k,
 \qquad x\longmapsto(x,x^{p^a},\ldots,x^{p^{a(b-1)}}).
\]
Thus its map is the direct sum of the rational representations
$N^{[aj]}$.  If a form pairs two summands, the corresponding summand
is expressed as the required dual.  A subfield inclusion is the same
natural matrix map, possibly precomposed with its specified field or
diagram automorphism.  On natural modules in the ranks in use, a
diagram automorphism gives $N$ or $N^*$; triality has already been
discarded with bounded rank.

\item In a twisted tensor orbit the map is
\[
 g\longmapsto
       \rho_0(F^{a_0}g)\otimes\cdots\otimes\rho_{b-1}(F^{a_{b-1}}g),
\]
where each $\rho_j$ is natural or dual.  This is a rational
representation and it is the defining matrix construction of the
finite subgroup.  Its natural module is
$\bigotimes_j N_j^{[a_j]}$, $N_j\in\{N,N^*\}$.
The orbit lengths, field degrees, and omitted tensor dimensions are
bounded by Proposition~\ref{grw:remaining-paths}.
\end{enumerate}

In characteristic two, two further classical maps must be interpreted
literally.  The nonsingular-vector stabilizer in a group of type
$D_{r+1}$ has connected reduced subgroup of type $B_r$; its action on
the $2r+2$-dimensional natural space has, by
Lemma~\ref{lift:nonsingular-line}, the invariant filtration
\[
 0<\langle v\rangle<v^\perp<W
\]
with factors $k,N_{B_r},k$.  Its simply connected cover therefore
supplies the required $B_r$-to-$D_{r+1}$ map.  The middle symplectic
action is an isogeny from this \emph{adjoint} $B_r$ stabilizer to
$\operatorname{Sp}_{2r}$.  It has no nontrivial
kernel on $k$-points, consistently with
Lemma~\ref{lift:nonsingular-line}; it is not a finite unipotent
radical.  In the following display, instead,
$\mathbf B_r=\operatorname{Spin}_{2r+1}$ and
$\mathbf C_r=\operatorname{Sp}_{2r}$ denote the \emph{simply
connected} groups.  The preceding vector stabilizer is the adjoint
quotient of $\mathbf B_r$; the kernels of its symplectic action and
of the simply connected maps are not being identified.  For the two
special isogenies
\[
 s:\mathbf B_r\longrightarrow\mathbf C_r,
 \qquad t:\mathbf C_r\longrightarrow\mathbf B_r,
 \qquad st=F_p,\quad ts=F_p,
\]
one has
\[
 s^*N_C=N_B,\qquad t^*N_B=N_C^{[1]}
\]
by \cite[author preprint, \S I.3.2]{DowdSin1996}.
If the inverse of $s$ on $\mathbb F_q$-points is needed, it is induced
by the rational map $tF_p^{f-1}$, since $q=p^f$.  No inverse
inseparable morphism is asserted.  These formulas preserve the
claimed description of simple constituents, including the trivial
factors of the full odd orthogonal module.

\smallskip\noindent
\emph{Rational maps for the nonlinear edges.}
For a defining-characteristic almost-simple edge, inflate its actual
finite irreducible representation along $\pi_i:U_i\twoheadrightarrow L_i$,
using the unique linear lift of Lemma~\ref{lift:finite-covers}
if the natural action was given
projectively.  Surjectivity preserves every invariant subspace and
every prescribed form.  A form similitude multiplier is a character
of the perfect group $U_i$, hence is trivial.
Only now apply Steinberg's restriction
Theorem~\cite[Theorem~43]{SteinbergNotes}: the resulting $U_i$-module
is the restriction of its $q_i$-restricted rational simple module,
with the specified field and diagram twist.  Here
$\sigma_i=F_{p^{f_i}}\gamma_i$ uses $f_i$ ordinary Frobenius digits,
including in the unitary case; the unitary natural scalar field
$\mathbb F_{p^{2f_i}}$ does not change this restriction parameter.
This yields a rational map to
$\operatorname{GL}(V)$ whose determinant is trivial: a connected
semisimple group has no nontrivial rational character.  Thus it maps
to a linear target, including the algebraic group underlying a
finite unitary target.  In a symplectic target or an orthogonal target
in odd characteristic, finite self-duality and the injective
highest-weight parametrization identify the invariant rational
bilinear form with the finite form, by Schur's Lemma.  Its symmetry
or alternating property is consequently the prescribed one.
For an orthogonal target in characteristic two the remaining
nonlinear module satisfies Proposition~\ref{lift:quadratic}, which
lifts its actual quadratic form.  The natural and special-isogeny
cases are the geometric maps already constructed.  The image is
connected and hence lies in the identity component of the isometry
group.

\smallskip\noindent
\emph{Compatibility through the fixed central covers.}
These constructions initially give maps into whichever central form
of the target is specified by its natural action.  A homomorphism
from a simply connected semisimple source lifts through a central
isogeny of semisimple algebraic groups; see
\cite[author version, Exercise~6.5.2(iii), p.~236]{Conrad2014}.  Thus they give maps
between simply connected covers as claimed.  Let
$\phi_i:\mathbf G_i\to\mathbf G_{i+1}$ be such a map and let
$\iota_i:L_i\hookrightarrow L_{i+1}$ be the actual finite inclusion.
By construction their induced maps to the target adjoint group
agree on $U_i$.  For $u\in U_i$, the target adjoint image of
$\phi_i(u)$ is fixed by $\sigma_{i+1}$, so
\[
 \delta(u)=\sigma_{i+1}(\phi_i(u))\phi_i(u)^{-1}
       \in Z(\mathbf G_{i+1})(k).
\]
Centrality gives $\delta(uv)=\delta(u)\delta(v)$.
Perfectness of $U_i$ forces $\delta=1$, so
$\phi_i(U_i)\leq U_{i+1}$, including when the algebraic centre is
nonreduced.  The two homomorphisms
$\pi_{i+1}\phi_i|_{U_i}$ and $\iota_i\pi_i$ to $L_{i+1}$ agree
modulo its centre.  Their discrepancy is consequently a
homomorphism $U_i\to Z(L_{i+1})$, and again is trivial.  Thus we have
the exact commutative diagram
\begin{equation}\label{lift:finite-diagram}
\begin{array}{ccc}
 U_i&\xrightarrow{\ \phi_i|_{U_i}\ }&U_{i+1}\\
 \pi_i\downarrow&&\downarrow\pi_{i+1}\\
 L_i&\xrightarrow{\ \iota_i\ }&L_{i+1}.
\end{array}
\end{equation}
This proves compatibility with the actual finite inclusions through
the fixed central-cover maps.  It gives no uniqueness assertion for
rational maps agreeing on finite points.
\end{proof}

\subsection{Composing a chosen path}

\begin{lemma}[Pathwise lifting]\label{lift:path-lifts}
Fix a path through retained edges with finite nested components
$L_0\leq\cdots\leq L_s$.  Choose an edge lift
$\phi_i:\mathbf G_i\to\mathbf G_{i+1}$ from
Proposition~\ref{lift:edge-lifts} for each edge, with the finite
central-cover identifications specified there.  Then
\[
 \Phi_{i,j}=\phi_{j-1}\cdots\phi_i\qquad(i<j)
\]
is a rational lift of the actual finite inclusion along that path.
If a rational $\mathbf G_j$-module is irreducible on $L_i$ under this
finite inclusion, it is irreducible on the algebraic image
$\Phi_{i,j}(\mathbf G_i)$.  A nontrivial such image is a closed
connected simple algebraic group; if its rank is strictly smaller
than the rank of $\mathbf G_j$, it is proper.
\end{lemma}

\begin{proof}
The finite compatibility statements for the edges compose, so the
first assertion follows by induction on $j-i$.  Every subspace
invariant under the algebraic image is invariant under the contained
finite image.  Finite irreducibility therefore implies algebraic
irreducibility, without a density assertion or a bound on the
highest weight of the composite representation.  The kernel of a
nontrivial map from a connected simple algebraic group is finite as
a group scheme: its connected reduced normal subgroup cannot be
the whole source.  Its image is closed and connected, and is simple
up to the allowed central form.  Inseparability can interchange the
types $B$ and $C$ in characteristic two, but does not change rank.
The last assertion follows from the rank inequality.
\end{proof}

\begin{remark}\label{lift:frobenius-periods}
The exponents in each chosen $\phi_i$ are retained when composing.
For example, on $\operatorname{SL}_n(p^f)$ the representations
$g\mapsto g\otimes F^a(g)$ and
$g\mapsto g\otimes F^{a+f}(g)$ agree, although their rational
representations need not agree.  Accordingly an independently
chosen shortcut map is never identified with a composed path map.
The only comparison needed later is their agreement on the actual
finite component, to which the irreducibility implication in
Lemma~\ref{lift:path-lifts} applies.  Also, a strict finite subfield
inclusion can have full algebraic image; properness in the later
application comes from rank growth.
\end{remark}

\begin{corollary}[Passing irreducibility backwards through geometry]
\label{lift:geometric-irreducibility}
Suppose that a retained geometric edge has source component $L$,
and that the target natural module restricts absolutely irreducibly
and nontrivially to a subgroup $H\leq L$.  Then the source natural
simple module restricts absolutely irreducibly to $H$.
\end{corollary}

\begin{proof}
Work over the algebraic closure.  Any proper submodule or nontrivial
composition filtration for the $L$-action would be an $H$-invariant
one.  Thus the target module is a single nontrivial composition
factor for $L$.  Proposition~\ref{lift:edge-lifts} expresses it as a
tensor product of twists of the source natural module or its dual.
If one factor had a nonzero proper $H$-invariant subspace, tensoring
that subspace with all the other factors would give a nonzero
proper invariant subspace of the product.  Hence each factor is
irreducible on $H$.  Taking a dual or applying a field Frobenius to
all matrix entries preserves irreducibility, so the source natural
module itself is irreducible on $H$.
\end{proof}
